% A.4.25 reserved by root.  New, separate source after checkpoint05.
\subsection{A finite transport square with simultaneous geometric fusion}

Retain the marked free surface group \(\Pi=\Pi_{1,9}\),
its quaternion kernel \(T\), and the quotients
\(Q_f,Q_3,Q_L,Q_{\rm out}=Q_f\times Q_L\) of A.4.24.
Let \(D_T\) be the normal closure in \(T\) of its thirty-six
oriented peripheral loops. Define
\[
 \begin{gathered}
 K_R=\gamma_3T\,T^3[T,D_T],\qquad R=T/K_R,
       \qquad E_R=\Pi/K_R,\\
 m=128\cdot3^{1172},\qquad
 Q_0=\Pi/\kappa_m(\Pi).
 \end{gathered}
 \tag{TF.1}
\]
Here \(\kappa_m\) is the characteristic diagonal kernel of
all homomorphisms into \(\mathfrak S_m\). The group \(Q_0\)
retains all the named peripheral elements. Let \(G_0\) be the
actual image in \(\operatorname{Out}(Q_0)\) of the
orientation-preserving nine-point mapping class group,
with its specified based lifts. Define the finite transport group
\[
 \mathcal A_0=
 \{F\in\operatorname{Aut}(Q_0):
       F(C_i)\sim C_i\ (1\le i\le9),\quad
       F G_0F^{-1}=G_0\}.
 \tag{TF.2}
\]
The conjugation in (TF.2) is the natural action of
\(\operatorname{Aut}(Q_0)\) on \(\operatorname{Out}(Q_0)\).
Thus (TF.2) specifies the multiplier-one peripheral transport
domain, together with its compatible action on the actual
geometric image. It does not identify that domain with the
geometric image.

\begin{result}{A.4.25}{Proved}{Simultaneous finite fusion and the ordinary descent square}
\label{thm:finite-transport-fusion}
The group \(Q_0\) has the actual peripheral-preserving
quotient maps
\[
 Q_0\longrightarrow Q_f\times Q_3
        \longrightarrow Q_f\times Q_L=Q_{\rm out},
 \qquad Q_0\longrightarrow E_R,
 \tag{TF.3}
\]
and \(|R|=3^{776}\), \(|E_R|=8\cdot3^{776}\).
The output in (TF.3) is preserved by the full geometric
group and the actual deck translations. The single \(Q_3\)
projection is used on the pure transport domain (TF.2).

Let \(G\) be the actual geometric image in
\(\operatorname{Out}(Q_{\rm out})\), and let
\(H\simeq C_3^2\) be the actual, labelled deck subgroup
of the three-division covering. For every \(F\in\mathcal A_0\),
there are actual pure geometric words \(\theta,w\) such that
\(F'=w^{-1}\theta^{-1}F\), at the output, fixes the
\(Q_L\) outer projection and the signed compact period
modulo twelve, and centralizes \(H\) pointwise. The word
\(w\) itself has trivial \(Q_L\) outer projection and
trivial signed compact period, and is in the quaternion kernel.

For the actual geometric normalizer there is a surjective
homomorphism
\[
        d:N_G(H)\longrightarrow J_c,
 \tag{TF.4}
\]
where \(J_c\) is the actual ordinary modular image on
\(\mathcal P_c\), of order \(1152\cdot3^4\).
The corrected transport defines an outer automorphism
\(\widetilde F'\) of \(\mathcal P_c\), and the exact square is
\[
 \begin{array}{ccc}
 N_G(H)&\xrightarrow{\ \operatorname{Ad}(F')\ }&N_G(H)\\
 d\downarrow&&\downarrow d\\
 J_c&\xrightarrow{\ \operatorname{Ad}(\widetilde F')\ }&J_c.
 \end{array}
 \tag{TF.5}
\]
In particular the corrected transport acts on the ordinary
image through this square. The normalizer-to-ordinary comparison
in (TF.4) is applied to the geometric elements \(n\in N_G(H)\).
\end{result}
\begin{proof}
\par\medskip\noindent\textbf{The finite cover form, with its normalization.}\quad
The index-eight cover \(T\) has genus nineteen and
thirty-six boundary components. Indeed the Euler
characteristic of \(\Pi_{1,9}\) is \(-9\), so the cover
has Euler characteristic \(-72\). Each original boundary
has quaternion inertia of order two and consequently four
lifts, giving thirty-six boundaries. Thus
\(2-2g-36=-72\), or \(g=19\).
Choose the usual oriented surface presentation for this cover,
with thirty-eight handle generators \(a_i,b_i\) and thirty-five
independent boundary generators \(\delta_j\), the last one
specified by the surface relation.

In \(R\), all boundary elements are central of order three.
Class-two, exponent-three collection gives independent
coordinates consisting of the seventy-three abelian
generators and the \(\binom{38}{2}=703\) commutators of
the handle generators. A model is the vector space
\(\mathbb F_3^{73}\oplus\bigwedge^2\mathbb F_3^{38}\)
with multiplication
\[
 (v,z)(v',z')=
 (v+v',z+z'+\tfrac12(\bar v\wedge\bar v')),
 \tag{TF.6}
\]
where the bar kills the thirty-five boundary coordinates.
Here \(\tfrac12=2\) in \(\mathbb F_3\).
This model attains the collection bound and proves
\(|R|=3^{73+703}=3^{776}\). The defining kernel is
normal in \(\Pi\): its deck action permutes the lifted
peripheral conjugacy classes. Hence \(E_R\) exists and
has the stated order.

Put \(V=H_1(T,\mathbb F_3)/\langle\delta_j\rangle\),
of dimension thirty-eight. The surface relation in \(R\)
is the exact central identity
\[
 \prod_{j=1}^{36}\delta_j=\omega^{-1},
 \qquad \omega=\sum_{i=1}^{19}[a_i,b_i]
                         \in\bigwedge^2V.
 \tag{TF.7}
\]
Because the boundary loops are central, this product is
independent of their ordering and of every change of their
based paths. A peripheral automorphism with common power
\(\lambda\) permutes these loops and sends each to its
\(\lambda\)-th power. Thus its induced compact matrix
\(M\) satisfies, with no undetermined scalar,
\[
                (\bigwedge^2M)\omega=\lambda\omega.
 \tag{TF.8}
\]
The inverse matrix of this nondegenerate alternating
bivector is the compact intersection form \(J_T\), and
(TF.8) is equivalent to \(M^tJ_TM=\lambda J_T\).

The central quaternion involution fixes every boundary
class, since it is the nonidentity element of the inertia
of each lifted boundary. Therefore the boundary radical
lies in its positive eigenspace. The negative eigenspace
of dimension thirty-six embeds into \(V\), and the
intersection form on it is nondegenerate: positive and
negative eigenspaces are perpendicular under the
quaternion-invariant form, and \(J_T\) is nondegenerate.
In the specified primitive multiplicity coordinates,
commutation with the quaternion matrices turns (TF.8)
into
\[
                         a^t F a=\lambda F.
 \tag{TF.9}
\]
Here is an explicit verification of the comparison to the
already retained matrix \(F\), including the scalar.
Let \(P:\mathbb F_3^{73}\to V\) be the quotient by the
thirty-five actual boundary homology vectors, let \(C_B\)
be the retained \(73\)-by-\(18\) primitive basis matrix,
and let \(D_i\) be the actual quaternion deck matrix.
The actual symmetric multiplicity form is
\[
             (PC_B)^tJ_T(PD_iC_B)=2F.
 \tag{TF.10}
\]
The identity is checked entry by entry by
\texttt{check\_main\_cover\_form.py}. That verifier rewrites
all thirty-six oriented boundary words in the actual
seventy-three-letter Schreier basis, forms their product
under (TF.6), and obtains a nondegenerate bivector of
rank thirty-eight. It checks preservation by all fifteen
geometric source matrices and verifies (TF.10).
The form in (TF.10) is symmetric because \(D_i\) is an
isometry with \(D_i^2=-I\); its nondegeneracy follows
also from the fact that \(D_i\) moves the primitive
quaternion line to an independent line. An automorphism
commuting with the quaternion action preserves \(C_B\)'s
image and commutes with \(D_i\). Applying (TF.8) to
(TF.10) therefore proves (TF.9) for that automorphism.
For the multiplier-one domain, \(a\) is an isometry.

\par\medskip\noindent\textbf{The characteristic input and all quotient maps.}\quad
The product of the orders of \(E_R\), the marked compact
abelian quotient modulo four, and \(Q_3\) is
\(8\cdot3^{776}\cdot16\cdot3^{396}=m\).
Their actual joint image is a quotient of \(\Pi\) of
order at most \(m\). Its left regular action, with fixed
letters added to degree \(m\), shows that its kernel
contains \(\kappa_m(\Pi)\). Thus all those actual maps
factor through \(Q_0\). The diagonal construction also
gives the explicit finite bound
\[
                       |Q_0|\le(m!)^{(m!)^{10}}.
 \tag{TF.11}
\]

We identify the relative kernels inside \(Q_0\) so that
their invariance is proved at this finite input.
The unique epimorphism-kernel to \(Q_8\) sending every
named peripheral element to \(-1\) recovers \(T\),
as in A.4.24: the two handle images must generate
\(Q_8\), and its twenty-four generating pairs are
one \(\operatorname{Aut}(Q_8)\)-orbit.
In the recovered subgroup \(T_0\) of \(Q_0\), take
the normal closure of the four conjugates of each
\(C_j^2\), one for each coset of its quaternion
inertia. This is the actual image \(D_0\) of \(D_T\).
A peripheral automorphism permutes these thirty-six
classes with the same common power, so preserves this
normal closure. Explicitly, after applying the
automorphism to a representative \(rC_j^2r^{-1}\),
write the new conjugating element as
\(t r' C_j^\epsilon\), using its coset in
\(Q_8/\langle-1\rangle\), with \(t\in T_0\).
The factor \(C_j^\epsilon\) centralizes \(C_j^2\),
leaving the specified \(T_0\)-conjugate of
\(r'C_j^2(r')^{-1}\). This also proves the permutation
assertion directly in the finite input. The quotient
\[
          T_0/(\gamma_3T_0\,T_0^3[T_0,D_0])
                       \ =\ R
 \tag{TF.12}
\]
is exactly the projected quotient, since
\(\kappa_m(\Pi)\subseteq K_R\).
Moreover \(K_R\subseteq T'T^3\), so the mod-three
abelianization of \(T_0\) is the full specified
seventy-three-dimensional homology of \(T\).
This proves finite invariance of \(E_R\) and the
cover-form calculation, without lifting an arbitrary
finite automorphism to an unspecified infinite group.

The negative eigenspace of the central quaternion
involution on the mod-three abelianization of \(T_0\)
recovers the faithful homology quotient. Joining it with
the marked compact abelian quotient modulo four gives
\(Q_f\), by the intrinsic construction in A.4.24.
The verbal output \(Q_L\) is characteristic; the
single stronger factor \(Q_3\) is invariant under the
pure peripheral domain (TF.2). These recover precisely
the maps (TF.3). We use the stronger input to impose
the peripheral equations, and perform deck conjugation
on the preserved output \(Q_f\times Q_L\).
Both output factors are characteristic in that product,
and its automorphism and outer automorphism groups
are their actual products, by A.4.24.

We also need the determinant of a transport's signed
compact period modulo four. The characteristic input
retains the compact quotient
\[
 V_4=F(A,B)/
 \bigl(\gamma_3F(A,B)\,(\gamma_2F(A,B))^4F(A,B)^8\bigr),
                 \qquad |V_4|=8^2\cdot4=256,
 \tag{TF.13}
\]
with the eight independent puncture generators killed.
It has unique forms \(B^jA^i[A,B]^k\), with
\(i,j\bmod8\), \(k\bmod4\); collection, or the model
\((i,j,k)(i',j',k')=(i+i',j+j',k+k'+ij')\), proves
its order and its exponent-eight relations.
Its defining kernel after killing those eight punctures
is verbal in the resulting two-generator quotient,
so the pure peripheral transport preserves it.
The ninth boundary is the inverse commutator, of exact
order four. Its multiplier-one preservation gives
\(\det C=1\pmod4\). A.4.24 gives \(\det C=1\pmod9\)
from the retained \(Q_3\) compact quotient.
In particular the signed period modulo twelve belongs
to \(\mathrm{SL}_2(\mathbb Z/12)\).
The factor (TF.13) is already retained by \(Q_0\), since
its order is at most \(m\), and its relative quotient
kernel is recovered by the just specified verbal operations.

\par\medskip\noindent\textbf{First geometric matching, including the quaternion residue.}\quad
Fix \(F\in\mathcal A_0\). A.4.24 supplies an actual
pure, orientation-preserving mapping class \(\theta_0\)
with the same \(Q_L\) outer action as \(F\).
Its compact period consequently agrees with that of
\(F\) modulo nine. The two actual handle words
\(U^9,L^9\) act identically on \(Q_L\), since its
named handle generators have order nine. Modulo four
their signed period matrices are exactly \(U,L\),
which generate \(\mathrm{SL}_2(\mathbb Z/4)\).
Multiplying \(\theta_0\) by an appropriate word in
these two corrections gives a pure geometric
\(\theta\) which still has the same \(Q_L\) outer
action and now has exactly the same signed period
modulo twelve as \(F\).

The outer automorphism of the specified quaternion
quotient is its handle matrix modulo two: all the
peripheral values are central \(-1\), so vanish in
\(Q_8/Z(Q_8)\). Thus \(F\) and \(\theta\) have
the same outer quaternion residue. Their remaining
quaternion discrepancy is inner. Choose an actual
element of \(\Pi\) mapping to its conjugating quaternion
element and modify the based lift of \(\theta\) by
this inner automorphism. This does not change its
\(Q_L\) outer action, any homology action, or any
ordinary period. It makes the two quaternion
automorphisms equal, as genuine automorphisms.
With this based choice put
\[
                            a_0=\theta^{-1}F.
 \tag{TF.14}
\]
It has identity quaternion action, identity signed
period modulo twelve, and identity \(Q_L\) outer
action. Its action on the multiplicity space is
an isometry \(a\in O(M,F)\), by (TF.9).

\par\medskip\noindent\textbf{Faithful fusion while every lower coordinate is fixed.}\quad
By A.4.20 choose \(c\in C_{O(M,F)}(H)\) with
determinant and spinor inverse to those of \(a\).
Then
\[
               k=ac\in\Omega_{18}^-(3),\qquad
               khk^{-1}=aha^{-1}\quad(h\in H).
 \tag{TF.15}
\]
The last equality preserves the individual labels,
not only the abstract subgroup.
A.4.23 supplies an actual pure quaternion-kernel word
\(w\) whose multiplicity image is \(k\), and whose
three lower coordinates (PC.2) are the identity.
The injection (JP.7) then gives identity outer action
on \(Q_L\) itself, and its signed period is the identity.
Set \(F'=w^{-1}a_0\).
Its multiplicity matrix is \(k^{-1}a=c^{-1}\), which
centralizes each labelled element of \(H\).
Its quaternion and signed-period actions are still
the identity, and its \(Q_L\) outer action is still
the identity.

This matrix assertion gives the claimed outer group
assertion on the faithful factor with its exact map.
In the split model
\(Q_f=A_f\rtimes G_2\), the base action of \(F'\)
is the identity: the joint quaternion and mod-four
coordinates determine \(G_2\), and both are now
fixed. The \(H\)-actions have base action the identity
and the specified faithful module action.
Equality of compatible base and module pairs gives
equality of their outer automorphisms, because the
only remaining ambiguity is the inner \(A_f\)-action
proved in A.4.19. Thus \(F'\) centralizes the labelled
\(H\) in \(\operatorname{Out}(Q_f)\).
It centralizes its other projection because its
\(Q_L\) outer action is the identity.
The actual product isomorphism on outer automorphisms
now proves centralization of the same diagonal \(H\)
in \(\operatorname{Out}(Q_{\rm out})\).

Since \(F\) normalizes \(G_0\) and \(\theta,w\) are
actual geometric words, the corrected transport still
normalizes the geometric image at the output.
Its conjugation consequently gives an automorphism
of \(N_G(H)\). This proves all fusion assertions
without placing the transport itself inside that
geometric normalizer.

\par\medskip\noindent\textbf{Descent of the actual geometric normalizer.}\quad
For \(n\in N_G(H)\), its two output coordinates
induce the same automorphism \(\nu\) of the labelled
\(H\), because both projections of this diagonal
subgroup are faithful. A.4.24 gives the actual
comparison \(\nu=J(C\bmod3)J\) with the compact
period in the primitive upstairs basis.
The period modulo four is retained by \(Q_f\),
and its modulo-three coordinate by \(Q_L\); together
they give the signed modulo-twelve coordinate.

We first define \(d\) before choosing any modular
representative. Choose a split representative
\((f,\beta_2)\) of the \(Q_f\) outer coordinate
on \(A_f\rtimes G_2\), modulo the specified inner
module ambiguity. Its normalizer action makes
\(f\) preserve the submodule
\(\mathfrak m_H^2A_f\); let \(\bar f\) denote
its induced map on \(A_c=A_f/\mathfrak m_H^2A_f\).
The defining compatible pair for \(d(n)\) is
\[
                  (\bar f,\ \beta_2\times\nu)
       \quad\hbox{on}\quad A_c\rtimes(G_2\times H).
 \tag{TF.15a}
\]
Compatibility on \(G_2\) is inherited from the
\(Q_f\) automorphism, and compatibility on \(H\)
is the literal normalizer identity verified below.
Formula (EX.3) therefore defines this outer
automorphism of \(\mathcal P_c\).
Changing the representative by an inner
\((x,g_2)\in A_f\rtimes G_2\) changes (TF.15a)
by the corresponding inner
\((\bar x,g_2,1)\) of \(\mathcal P_c\).
Changes of the split coordinates have the same
inner-module ambiguity. Hence \(d\) is independent
of all these choices. Composition of the compatible
pairs proves that it is a homomorphism.
The following geometric comparison determines
its image.

First remove the modulo-two residue by an actual
ordinary modular lift, and remove the remaining
inner quaternion residue by an actual based inner
\(\Pi\)-correction. For the resulting geometric
element the multiplicity matrix is in
\(\mathrm{SO}_{18}^-(3)\), by A.4.15, and satisfies
the spinor/period identity \(\theta(g)=\eta(q)\).
The latter is checked on all 360 actual
quaternion-kernel generators, with their signed
periods; multiplicativity proves it on the whole
geometric kernel. Thus its first-order action and
signed period lie in precisely (NJ.7), whose image
is the actual ordinary one by A.4.21.

Here the normalizer condition on the module is
literal, even though the initial equality was outer.
In \(Q_f=A_f\rtimes G_2\), both conjugated deck
actions have identity base action. An inner difference
therefore comes from \(Z(G_2)\), whose image on
\(A_f\) is \(I\) or \(-I\). The negative choice
would turn an order-three matrix into one of order
six, so it cannot occur. Hence
\(g h g^{-1}=h_{\nu(h)}\) holds on the module itself.

The kernel \(\mathfrak m_H^2M\) is invariant under
this normalizer by the exact formula (QF.K4).
The induced action on the twelve-dimensional faithful
quotient, together with the specified signed base map
on \(G_2\times H\), is exactly (TF.15a).
The complete comparison and its inner kernel of
order eighteen are A.4.22. Restoring the six
modulo-two residues proves that its image is contained
in the full ordinary image \(J_c\).
Conversely the actual downstairs modular maps
\(u,l\) lift through the three-division cover and
normalize \(H\), so the image contains all of \(J_c\).
This proves (TF.4), including surjectivity and its
agreement with ordinary descent.

\par\medskip\noindent\textbf{The transport square.}\quad
The corrected transport has identity quaternion
and signed-period action and centralizes \(H\) on
the faithful module. It therefore preserves
\(\mathfrak m_H^2M\) and induces an actual map on
the compact twelve-dimensional quotient. Extend this
map over \(G_2\times H\) by the identity.
The module map commutes with both factors, so the
split-extension formula (EX.3) gives an automorphism
\(\widetilde F'\) of \(\mathcal P_c\).
Its outer class is independent of the chosen splitting:
the explicitly retained ambiguity is inner
conjugation by the twelve-dimensional module.

For every geometric \(n\in N_G(H)\), composing
these induced module and base maps in either order
gives
\[
 d(F'n(F')^{-1})=
       \widetilde F'\,d(n)\,(\widetilde F')^{-1}
                         \quad\hbox{in }\operatorname{Out}(\mathcal P_c).
 \tag{TF.16}
\]
The reconstruction theorem identifies any difference
of compatible representatives as its specified inner
module conjugation, which vanishes in this outer group.
This proves the equality rather than merely equality
of the two linear matrices. Surjectivity of \(d\)
now makes the lower conjugation an automorphism of
\(J_c\), and proves the square (TF.5).
The determinant and spinor/period hypotheses used in
this argument were those of the actual geometric
elements \(n\); the transport supplies the horizontal
automorphisms in the square.
\end{proof}
